\section{Giới thiệu}

\subsection{Bối cảnh và Tính cấp thiết của đề tài}

Nông nghiệp là ngành kinh tế đóng vai trò trụ cột tại Việt Nam, đặc biệt là các mô hình canh tác cây trồng ngoài trời (Open-field Farm) có giá trị kinh tế cao như cây ăn trái (sầu riêng, thanh long, bưởi tại Đồng bằng sông Cửu Long và Đông Nam Bộ) và cây công nghiệp (cà phê, hồ tiêu tại Tây Nguyên). Tuy nhiên, phương thức canh tác truyền thống trên diện tích đất mở hiện đang đối mặt với những thách thức và rủi ro rất lớn:
\begin{itemize}
    \item \textbf{Khí hậu khắc nghiệt, hạn hán và sốc nhiệt mùa khô:} Hiện tượng biến đổi khí hậu gay gắt dẫn đến các đợt nắng nóng cực đoan (nhiệt độ ngoài trời vượt $38 - 40^\circ\text{C}$) kết hợp độ ẩm không khí xuống thấp, làm gia tăng tốc độ bốc thoát hơi nước của cây trồng, gây cháy đọt, rụng hoa và kiệt sức bộ rễ nếu không được phun sương làm mát tầng tán hoặc cấp nước kịp thời.
    \item \textbf{Tưới tiêu và điều tiết ánh sáng theo cảm tính:} Nông dân thường tưới xả tràn hoặc canh giờ tưới dựa trên mắt nhìn bề mặt đất, dẫn đến lãng phí nguồn nước ngầm quý giá và thất thoát chi phí điện năng. Tương tự, ở các vùng chuyên canh hoa cúc hoặc thanh long, việc canh thời gian chong đèn quang kỳ để kích thích ra hoa trái vụ còn mang tính thủ công, thiếu tự động hóa.
    \item \textbf{Khan hiếm nhân công và khó khăn quản lý phân tán:} Với quy mô nông trại rộng nhiều hecta chia thành nhiều lô/khu đất (Zone/Plot), việc điều động nhân công đi kiểm tra thực địa, đóng mở từng van nước hoặc bật từng cầu dao điện phân tán ngoài đồng tốn rất nhiều thời gian, chi phí và khó ứng phó kịp khi máy bơm gặp sự cố.
    \item \textbf{Thiếu dữ liệu lịch sử phục vụ ra quyết định canh tác:} Mô hình truyền thống hoàn toàn thiếu vắng cơ chế ghi nhận nhật ký chuỗi thời gian về độ ẩm đất và vi khí hậu, khiến chủ nông trại không thể đối chiếu lượng nước tưới tiêu thụ với năng suất thu hoạch của từng lô đất qua các mùa vụ để tối ưu hóa chi phí.
\end{itemize}

Trước thực trạng đó, việc ứng dụng công nghệ Internet vạn vật (\gls{iot}) kết hợp với \gls{httt} quản lý tập trung, mô hình \textbf{Nông trại thông minh (Smart Farm)}, trở thành giải pháp thiết thực giúp tự động hóa khâu giám sát vi khí hậu, tối ưu hóa tài nguyên nước tưới và nâng cao năng lực quản trị nông nghiệp bền vững tại Việt Nam.

\subsection{Mục tiêu và Phạm vi của dự án}

\subsubsection{Mục tiêu tổng quát}
Mục tiêu cốt lõi của đề tài là xây dựng một hệ thống hoàn chỉnh từ tầng phần cứng cảm biến, tầng truyền thông mạng, cổng kết nối \gls{iot} Gateway đến tầng ứng dụng quản lý trung tâm, hỗ trợ người nông dân và chủ trang trại quản lý, giám sát và vận hành thiết bị nông trại một cách tự động, chính xác và hiệu quả.

\subsubsection{Mục tiêu cụ thể theo Module}
\begin{enumerate}[leftmargin=*]
    \item \textbf{Module 1 (Thu thập dữ liệu đầu vào):} Tự động đo đạc liên tục các thông số môi trường vi khí hậu quan trọng (độ ẩm đất, nhiệt độ, độ ẩm không khí, cường độ ánh sáng) tại các khu đất canh tác.
    \item \textbf{Module 2 (Kiểm tra dữ liệu \& Phát hiện vượt ngưỡng):} Xử lý dữ liệu thu thập theo thời gian thực, đối chiếu với dải ngưỡng an toàn tương ứng của từng loại cây trồng và phát hiện kịp thời các điều kiện môi trường bất thường.
    \item \textbf{Module 3 (Điều khiển thiết bị đầu ra):} Hỗ trợ song song hai cơ chế: điều khiển tự động (dựa trên quy tắc ngưỡng và lịch trình thời gian) và điều khiển thủ công từ xa đối với các cơ cấu chấp hành (máy bơm tưới nước, hệ thống phun sương hạ nhiệt, dàn đèn chong quang kỳ, van điện từ).
    \item \textbf{Module 4 (Nhật ký hoạt động ứng dụng - Audit Log):} Ghi nhận độc lập mọi biến động hệ thống và hành vi can thiệp của người dùng vào cơ sở dữ liệu nội bộ (không sử dụng log mặc định của IoT Server) nhằm phục vụ thanh tra, kiểm toán vận hành.
    \item \textbf{Module 5 (Data Dashboard \& Phân tích):} Xây dựng giao diện Dashboard trực quan hóa trạng thái toàn hệ thống theo thời gian thực và cung cấp công cụ báo cáo phân tích đa chiều, hỗ trợ trích xuất dữ liệu mùa vụ.
\end{enumerate}

\subsubsection{Phạm vi hệ thống}
Hệ thống được thiết kế theo cấu trúc phân cấp linh hoạt: Một nông trại được chia thành nhiều \textbf{Khu đất / Lô vườn canh tác (Zone/Plot)}. Mỗi khu đất được gán với một \textbf{Loại cây trồng (Crop Type)} cụ thể (ví dụ: lô trồng sầu riêng, vườn thanh long chong đèn, khu trồng cà phê, vườn cây có múi). Mỗi loại cây trồng sẽ sở hữu bộ cấu hình riêng biệt về dải ngưỡng sinh trưởng an toàn (ngưỡng độ ẩm đất, ngưỡng nhiệt độ chống sốc nhiệt), lịch trình tưới tiêu định kỳ và danh mục thiết bị phụ trách (bơm tưới gốc, béc phun sương làm mát tầng lá, dàn đèn quang kỳ).
